% GENERATED FILE. Edit canonical lessons, then run npm run generate:books.
% canonical-source-hash: fnv1a64:c433773ca201fe61
% canonical-lessons: MW-C78-din, MW-C78-rat, MW-C78-savero, MW-C78-mahino, MW-C78-baras

\chapter{Day, Night, Morning, Month, Year}
\label{ch:mw-day-night-morning-month-year}
\hlchaptermodality{78}

\begin{chapteropening}
\textbf{By the end of this chapter:} \emph{I can say a day, night, morning, a month and a year.}

Complete the last lesson of chapter 78: I can say a day, night, morning, a month and a year.
\end{chapteropening}

\section[\texorpdfstring{\mw{दिन} (din)}{दिन (din)}]{\mw{दिन} (din) — a day}

\label{lesson:MW-C78-din}



\begin{quote}
Before the new one: say the Marwadi for brown, then the Marwadi for sweet.
\end{quote}

\subsection*{You'll want to know: \mw{दिन}}
\textbf{\mw{दिन}} — \emph{din} — \textquotedblleft{}a day\textquotedblright{}.

\begin{cousinweb}
\textbf{\mw{दिन}} (\emph{din}) is Sanskrit \emph{dina}, \textquotedblleft{}day\textquotedblright{}. CDIAL lists Hindi and Gujarati \emph{din} under it and thinks the Gujarati and Marathi forms were taken back from Sanskrit (CDIAL 6328); it gives no Marwari form.
\end{cousinweb}

\begin{grammarlens}[title={What you've built}]
One more word for this chapter.
\end{grammarlens}

\subsection*{Your turn}
\begin{itemize}
\raggedright
  \item \emph{Say it:} \emph{din}
  \item \emph{Say it:} \emph{din}, once more
  \item \emph{Say it:} say \emph{din}
  \item \emph{Recall:} say the Marwadi for brown, then the Marwadi for sweet, then say \emph{din} again
\end{itemize}

\begin{tcolorbox}[breakable,colback=teal!4,colframe=teal!35!black,title={Before you move on}]
What does \mw{दिन} mean? (\textquotedblleft{}A day\textquotedblright{}.) Say it once more.
\end{tcolorbox}
\section[\texorpdfstring{\mw{रात} (rāt)}{रात (rāt)}]{\mw{रात} (rāt) — night}

\label{lesson:MW-C78-rat}



\begin{quote}
Before the new one: say the Marwadi for sweet, then the Marwadi for a day.
\end{quote}

\subsection*{You'll want to know: \mw{रात}}
\textbf{\mw{रात}} — \emph{rāt} — \textquotedblleft{}night\textquotedblright{}.

\begin{cousinweb}
\textbf{\mw{रात}} (\emph{rāt}) is a tadbhava, a word inherited from Sanskrit and reshaped by sound change. It comes from \emph{rātrī}, \textquotedblleft{}night\textquotedblright{}. CDIAL lists Old Marwari \emph{rāti} under it (CDIAL 10702).
\end{cousinweb}

\begin{grammarlens}[title={What you've built}]
One more word for this chapter.
\end{grammarlens}

\subsection*{Your turn}
\begin{itemize}
\raggedright
  \item \emph{Say it:} \emph{rāt}
  \item \emph{Say it:} \emph{rāt}, once more
  \item \emph{Say it:} say \emph{rāt}
  \item \emph{Recall:} say the Marwadi for sweet, then the Marwadi for a day, then say \emph{rāt} again
\end{itemize}

\begin{tcolorbox}[breakable,colback=teal!4,colframe=teal!35!black,title={Before you move on}]
What does \mw{रात} mean? (\textquotedblleft{}Night\textquotedblright{}.) Say it once more.
\end{tcolorbox}
\section[\texorpdfstring{\mw{सवेरो} (savero)}{सवेरो (savero)}]{\mw{सवेरो} (savero) — morning}

\label{lesson:MW-C78-savero}



\begin{quote}
Before the new one: say the Marwadi for a day, then the Marwadi for night.
\end{quote}

\subsection*{You'll want to know: \mw{सवेरो}}
\textbf{\mw{सवेरो}} — \emph{savero} — \textquotedblleft{}morning\textquotedblright{}.

\begin{cousinweb}
\textbf{\mw{सवेरो}} (\emph{savero}) is a tadbhava, a word inherited and reshaped by sound change. It comes from \emph{*savela}, \textquotedblleft{}in time, early\textquotedblright{}, a form CDIAL reconstructs from its descendants. CDIAL gives no Marwari form in this entry (CDIAL 13291); the match is through Hindi \emph{sawerā} and Gujarati \emph{saverā}. It is built on \emph{sa-}, \textquotedblleft{}with\textquotedblright{}, and \emph{velā}, \textquotedblleft{}time\textquotedblright{}.
\end{cousinweb}

\begin{grammarlens}[title={What you've built}]
One more word for this chapter.
\end{grammarlens}

\subsection*{Your turn}
\begin{itemize}
\raggedright
  \item \emph{Say it:} \emph{savero}
  \item \emph{Say it:} \emph{savero}, once more
  \item \emph{Say it:} say \emph{savero}
  \item \emph{Recall:} say the Marwadi for a day, then the Marwadi for night, then say \emph{savero} again
\end{itemize}

\begin{tcolorbox}[breakable,colback=teal!4,colframe=teal!35!black,title={Before you move on}]
What does \mw{सवेरो} mean? (\textquotedblleft{}Morning\textquotedblright{}.) Say it once more.
\end{tcolorbox}
\section[\texorpdfstring{\mw{महीनो} (mahīno)}{महीनो (mahīno)}]{\mw{महीनो} (mahīno) — a month}

\label{lesson:MW-C78-mahino}



\begin{quote}
Before the new one: say the Marwadi for night, then the Marwadi for morning.
\end{quote}

\subsection*{You'll want to know: \mw{महीनो}}
\textbf{\mw{महीनो}} — \emph{mahīno} — \textquotedblleft{}a month\textquotedblright{}.

\begin{cousinweb}
\textbf{\mw{महीनो}} (\emph{mahīno}) is a loan, usually traced to Persian \emph{māhīna}, \textquotedblleft{}monthly\textquotedblright{}. It is built on Persian \emph{māh}, \textquotedblleft{}moon, month\textquotedblright{}.
\end{cousinweb}

\begin{grammarlens}[title={What you've built}]
One more word for this chapter.
\end{grammarlens}

\subsection*{Your turn}
\begin{itemize}
\raggedright
  \item \emph{Say it:} \emph{mahīno}
  \item \emph{Say it:} \emph{mahīno}, once more
  \item \emph{Say it:} say \emph{mahīno}
  \item \emph{Recall:} say the Marwadi for night, then the Marwadi for morning, then say \emph{mahīno} again
\end{itemize}

\begin{tcolorbox}[breakable,colback=teal!4,colframe=teal!35!black,title={Before you move on}]
What does \mw{महीनो} mean? (\textquotedblleft{}A month\textquotedblright{}.) Say it once more.
\end{tcolorbox}
\section[\texorpdfstring{\mw{बरस} (baras)}{बरस (baras)}]{\mw{बरस} (baras) — a year}

\label{lesson:MW-C78-baras}



\begin{quote}
Before the new one: say the Marwadi for morning, then the Marwadi for a month.
\end{quote}

\subsection*{You'll want to know: \mw{बरस}}
\textbf{\mw{बरस}} — \emph{baras} — \textquotedblleft{}a year\textquotedblright{}.

\begin{cousinweb}
\textbf{\mw{बरस}} (\emph{baras}) is a tadbhava, a word inherited from Sanskrit and reshaped by sound change. It comes from \emph{varṣa}, \textquotedblleft{}rain; year\textquotedblright{}. CDIAL lists Old Marwari \emph{barasa} under it (CDIAL 11392). The sense moved from \textquotedblleft{}rain, the rainy season\textquotedblright{} to \textquotedblleft{}year\textquotedblright{}.
\end{cousinweb}

\begin{grammarlens}[title={What you've built}]
One more word for this chapter.
\end{grammarlens}

\subsection*{Your turn}
\begin{itemize}
\raggedright
  \item \emph{Say it:} \emph{baras}
  \item \emph{Say it:} \emph{baras}, once more
  \item \emph{Say it:} say \emph{baras}
  \item \emph{Recall:} say the Marwadi for morning, then the Marwadi for a month, then say \emph{baras} again
\end{itemize}

\begin{tcolorbox}[breakable,colback=teal!4,colframe=teal!35!black,title={Before you move on}]
What does \mw{बरस} mean? (\textquotedblleft{}A year\textquotedblright{}.) Say it once more.
\end{tcolorbox}
